
\section{2. Related Work}

\subsection{2.1 Sparse Attention Activations}

\textbf{Sparsemax.} Martins \& Astudillo (2016) introduce sparsemax, the $\alpha = 2$ 
case of the entmax family. Sparsemax projects scores onto the simplex, 
producing the sparsest possible distribution consistent with the scores. It 
has sparse outputs but is non-smooth at the support boundary, which can slow 
training.

\textbf{α-entmax.} Peters et al. (2019) generalize sparsemax to a one-parameter 
family. $\alpha = 1.5$ is the most commonly used value, balancing sparsity 
and smoothness. They provide forward and backward algorithms via bisection.

\textbf{Adaptive entmax.} Correia et al. (2019) extend α-entmax to allow 
$\alpha$ to be \textbf{learned per head}. We use fixed $\alpha = 1.5$ as default 
for simplicity, but the mechanism supports learned $\alpha$.

\textbf{Fused entmax.} Recent implementations (Zamora et al., 2024) provide 
fused CUDA kernels for entmax. Our implementation is a reference bisection.

\subsection{2.2 Local and Sliding-Window Attention}

\textbf{Longformer.} Beltagy et al. (2020) introduce a sliding-window attention 
of width $W$, giving $\mathcal{O}(L \cdot W)$ complexity. AdaSplash uses the 
same window but with α-entmax activations.

\textbf{BigBird.} Zaheer et al. (2020) combine local, random, and global attention. 
Their window is fixed-size; AdaSplash uses a dynamic splash radius.

\textbf{Blockwise attention.} Qiu et al. (2020) use block-sparse attention with 
fixed block size. Similar idea, different primitive.

\textbf{Sparse Transformer.} Child et al. (2019) use strided and fixed attention 
patterns. Orthogonal to our approach.

\subsection{2.3 Attention Sinks}

\textbf{StreamingLLM.} Xiao et al. (2024) discover that preserving the first 
few tokens ("attention sinks") is critical for stability in streaming 
inference. AdaSplash preserves the first 4 tokens by default.

\textbf{H2O.} Zhang et al. (2023) use an attention-sink + heavy-hitter strategy 
for KV-cache eviction. AdaSplash uses attention sinks for a different purpose: 
to preserve stability under local windowing.

\subsection{2.4 Parameter-Efficient Fine-Tuning}

\textbf{LoRA.} Hu et al. (2021) add trainable rank-$r_L$ updates $B A$ with 
$B = 0$. AdaSplash uses standard LoRA for task adaptation.

\textbf{AdaLoRA.} Zhang et al. (2023) adaptively allocate LoRA rank per layer. 
Orthogonal extension.

\textbf{PiSSA.} Meng et al. (2024) initialize LoRA adapters with principal 
components. Compatible with AdaSplash.

\subsection{2.5 Efficient Attention Mechanisms}

\begin{table}[H]
\centering
\begin{tabular}{lllll}
\toprule
\textbf{Method} & \textbf{Type} & \textbf{Complexity} & \textbf{Sparse?} & \textbf{Sinks?} \\
\midrule
Softmax attention & Dense & $\mathcal{O}(L^2)$ & No & No \\
Longformer & Local & $\mathcal{O}(L W)$ & No & No \\
BigBird & Local + random & $\mathcal{O}(L W)$ & No & No \\
Sparsemax & Activation & $\mathcal{O}(L^2)$ & Yes & No \\
Entmax & Activation & $\mathcal{O}(L^2)$ & Yes & No \\
StreamingLLM & Sink + window & $\mathcal{O}(L W)$ & No & Yes \\
**AdaSplash** & **Sink + window + entmax** & **$\mathcal{O}(L W)$** & **Yes** & **Yes** \\
\bottomrule
\end{tabular}
\end{table}

\subsection{2.6 Position of Our Work}

AdaSplash is unique in combining:
1. \textbf{α-entmax sparsity} at the activation level (not mask level).
2. \textbf{Local splash radius} with sink preservation.
3. \textbf{Multi-architecture support} including GQA and RoPE.
4. \textbf{Frozen base + LoRA} for parameter-efficient adaptation.

The closest prior work is the \textbf{combination} of Longformer (local window) 
with α-entmax (sparsity). To our knowledge, no published work evaluates this 
specific combination on small pretrained models with LoRA recovery.

\subsection{2.7 Important Distinction}

Unlike low-rank compression methods (A3, AASVD, Adaptive SVD), AdaSplash:
\begin{itemize}
\item Does \textbf{not} reduce parameter count.
\end{itemize}
\begin{itemize}
\item Does \textbf{not} compress weights.
\end{itemize}
\begin{itemize}
\item \textbf{Increases} parameter count by ~0.88\% via LoRA.
\end{itemize}

Its value is in \textbf{sparsity} and \textbf{interpretability}, not size.